% Combined TeX source bundle
% Title: Sunzi, Sunzi Suanjing, classical-modern Chinese
% Note: constituent files are preserved below in sources/. Some are standalone
% documents with their own preambles; this combined file is for inspection,
% comparison, and reuse rather than guaranteed one-command compilation.


% ============================================================
% Source: translations/zh/chinese/sunzi-suanjing_classical-modern_bilingual.tex
% ============================================================

%% ========================================================================
%% Sunzi Suanjing (孙子算经) — Classical Chinese ⟷ Modern Chinese Bilingual Edition
%% 孙子算经 — 文言文/白话文对照版
%% LEFT:  Original Classical Chinese (文言文)
%% RIGHT: Modern Chinese Vernacular (白话文)
%% ========================================================================
\documentclass[11pt,a4paper]{article}

%% -------- Core Packages --------
\usepackage{fontspec}
\usepackage{polyglossia}
\usepackage{amsmath,amssymb,amsthm}
\usepackage{geometry}
\usepackage{graphicx}
\usepackage{xcolor}
\usepackage{titlesec}
\usepackage{fancyhdr}
\usepackage{enumitem}
\usepackage{array,booktabs,longtable,multirow}
\usepackage{hyperref}
\usepackage{tocloft}
\usepackage{indentfirst}
\usepackage{tcolorbox}
\usepackage{etoolbox}
\usepackage{fancyvrb}
\usepackage{url}
\usepackage{paracol}
\usepackage{xeCJK}

%% -------- Page Geometry --------
\geometry{
  paperwidth=210mm,paperheight=297mm,
  left=16mm,right=16mm,top=24mm,bottom=24mm,
  headheight=14pt,footskip=28pt
}

%% -------- Language Setup --------
\setmainlanguage{chinese}

%% -------- Font Configuration --------
\setmainfont{Linux Libertine O}[
  Extension=.otf,
  UprightFont=*-Regular,
  BoldFont=*-Bold,
  ItalicFont=*-Italic,
  BoldItalicFont=*-BoldItalic,
  Renderer=Harfbuzz
]
\setsansfont{Linux Biolinum O}[
  Extension=.otf,
  UprightFont=*-Regular,
  BoldFont=*-Bold,
  ItalicFont=*-Italic,
  Renderer=Harfbuzz
]
\setmonofont{DejaVu Sans Mono}[Scale=MatchLowercase]

%% -------- xeCJK Setup --------
\setCJKmainfont{NotoSansCJKsc-Regular.otf}[
  Path=/mnt/agents/fonts/noto-sc-extracted/,
  BoldFont=NotoSansCJKsc-Bold.otf
]
\setCJKsansfont{NotoSansCJKsc-Regular.otf}[
  Path=/mnt/agents/fonts/noto-sc-extracted/,
  BoldFont=NotoSansCJKsc-Bold.otf
]
\setCJKmonofont{NotoSansCJKsc-Regular.otf}[
  Path=/mnt/agents/fonts/noto-sc-extracted/,
  BoldFont=NotoSansCJKsc-Bold.otf
]

%% -------- Colors --------
\definecolor{titlecolor}{RGB}{45,55,72}
\definecolor{sectioncolor}{RGB}{55,65,81}
\definecolor{notecolor}{RGB}{120,53,15}
\definecolor{rulecolor}{RGB}{180,160,140}
\definecolor{chinesebg}{RGB}{255,252,245}
\definecolor{problembg}{RGB}{250,248,244}
\definecolor{answerbg}{RGB}{245,250,245}
\definecolor{methodbg}{RGB}{245,248,252}
\definecolor{crtgold}{RGB}{139,90,0}
\definecolor{classicalbg}{RGB}{255,253,248}
\definecolor{modernbg}{RGB}{248,252,255}
\definecolor{columnsepbg}{RGB}{220,210,195}

%% -------- Column separation --------
\columnsep=16pt
\setlength{\columnseprule}{0.6pt}
\def\columnseprulecolor{\color{columnsepbg}}

%% -------- Section Formatting --------
\titleformat{\section}
  {\normalfont\LARGE\bfseries\color{titlecolor}}
  {\thesection}{1em}{}
\titleformat{\subsection}
  {\normalfont\large\bfseries\color{sectioncolor}}
  {\thesubsection}{1em}{}
\titleformat{\subsubsection}
  {\normalfont\normalsize\bfseries\color{sectioncolor}}
  {\thesubsubsection}{1em}{}

%% -------- Header/Footer --------
\pagestyle{fancy}
\fancyhf{}
\fancyhead[L]{\small\textsc{孙子算经 · 文言文/白话文对照版}}
\fancyhead[R]{\small\thepage}
\renewcommand{\headrulewidth}{0.4pt}
\renewcommand{\footrulewidth}{0pt}

%% -------- Custom Environments --------
\newtcolorbox{classicalbox}{
  colback=classicalbg,colframe=rulecolor!60,boxrule=0.4pt,
  left=6pt,right=6pt,top=4pt,bottom=4pt,arc=2mm
}
\newtcolorbox{modernbox}{
  colback=modernbg,colframe=rulecolor!60,boxrule=0.4pt,
  left=6pt,right=6pt,top=4pt,bottom=4pt,arc=2mm
}
\newtcolorbox{problemclassical}{
  colback=problembg,colframe=rulecolor,boxrule=0.5pt,
  left=6pt,right=6pt,top=4pt,bottom=4pt,fonttitle=\bfseries,title=问
}
\newtcolorbox{problemmodern}{
  colback=problembg,colframe=rulecolor,boxrule=0.5pt,
  left=6pt,right=6pt,top=4pt,bottom=4pt,fonttitle=\bfseries,title=题目
}
\newtcolorbox{answerclassical}{
  colback=answerbg,colframe=rulecolor!70!green,boxrule=0.5pt,
  left=6pt,right=6pt,top=3pt,bottom=3pt,fonttitle=\bfseries,title=答
}
\newtcolorbox{answermodern}{
  colback=answerbg,colframe=rulecolor!70!green,boxrule=0.5pt,
  left=6pt,right=6pt,top=3pt,bottom=3pt,fonttitle=\bfseries,title=答案
}
\newtcolorbox{methodclassical}{
  colback=methodbg,colframe=rulecolor!70!blue,boxrule=0.5pt,
  left=6pt,right=6pt,top=3pt,bottom=3pt,fonttitle=\bfseries,title=术
}
\newtcolorbox{methodmodern}{
  colback=methodbg,colframe=rulecolor!70!blue,boxrule=0.5pt,
  left=6pt,right=6pt,top=3pt,bottom=3pt,fonttitle=\bfseries,title=解法
}

%% -------- Custom Commands --------
\newcommand{\longline}{\noindent\rule{\textwidth}{0.4pt}}
\newcommand{\classtitle}[1]{\textbf{\color{titlecolor}#1}}
\newcommand{\modtitle}[1]{\textbf{\color{sectioncolor}#1}}
\newcommand{\bigcolswitch}{\switchcolumn[0]*}

%% Bilingual pair command
\newcommand{\bilingualsection}[2]{%
  \section*{#1}
  \addcontentsline{toc}{section}{#1}
}

\newcommand{\bilingualsubsection}[2]{%
  \subsection*{#1}
}

%% -------- Hyperref Setup --------
\hypersetup{
  colorlinks=true,
  linkcolor=sectioncolor,
  citecolor=sectioncolor,
  urlcolor=sectioncolor,
  pdfencoding=auto,
  pdftitle={孙子算经 — 文言文/白话文对照版},
  pdfauthor={孙子（约公元3—5世纪）}
}

%% -------- Sync columns --------
\synccounter{section}
\synccounter{subsection}
\synccounter{subsubsection}

%% ========================================================================

\begin{document}

%% =======================================================================
%% TITLE PAGE
%% =======================================================================
\begin{titlepage}
\begin{center}
\vspace*{1.5cm}

{\fontsize{12}{14}\selectfont\textsc{中国古典数学典籍}}\\[0.5cm]

\longline\\[1cm]

{\fontsize{32}{38}\selectfont\bfseries\color{titlecolor}
孙子算经}\\[0.6cm]

{\fontsize{20}{24}\selectfont\itshape Sunzi Suanjing}\\[0.4cm]

{\fontsize{14}{18}\selectfont Master Sun's Mathematical Manual}\\[1cm]

\longline\\[1cm]

\begin{tcolorbox}[colback=chinesebg,colframe=rulecolor,boxrule=1pt,
  width=0.9\textwidth,arc=3mm]
\centering
\large
\textbf{文言文 / 白话文 对照版}\\[4pt]
\normalsize
Classical Chinese / Modern Chinese Bilingual Edition
\end{tcolorbox}

\vspace{1cm}

{\large 原题：\textbf{孙子}（约公元3—5世纪）}\\[0.3cm]
{\normalsize 白话译文：现代汉语翻译}\\[1.5cm]

\begin{tcolorbox}[colback=chinesebg,colframe=rulecolor,boxrule=0.8pt,
  width=0.85\textwidth,arc=2mm]
\centering
\small
\textbf{《算经十书》之一}\\[3pt]
唐代（618—907年）列入官学教材\\[3pt]
\textbf{载有世界最早的中国剩余定理（孙子定理）}
\end{tcolorbox}

\vfill

{\small 用 \LaTeX + XeLaTeX + xeCJK 排版}\\[0.3cm]
{\small\today}

\end{center}
\end{titlepage}

%% =======================================================================
%% TABLE OF CONTENTS
%% =======================================================================
\tableofcontents
\newpage

%% =======================================================================
%% INTRODUCTION
%% =======================================================================
\begin{paracol}{2}

\begin{classicalbox}
\section*{简介}
\addcontentsline{toc}{section}{简介}

《孙子算经》是中国现存最早的数学著作之一，成书于公元3至5世纪。唐代将其列为《算经十书》之一。作者孙子的生平已不可考。

全书共分三卷：
\begin{itemize}[leftmargin=1.5em,topsep=2pt,itemsep=1pt]
  \item \textbf{卷上}：度量衡制度、大数记法、算筹记数、乘除法则、九九乘法表
  \item \textbf{卷中}：分数运算、面积体积计算、粮谷折算
  \item \textbf{卷下}：三十六道应用题，包括著名的「物不知数」题——中国剩余定理的最早记载
\end{itemize}
\end{classicalbox}

\switchcolumn

\begin{modernbox}
\section*{Introduction}

The \textit{Sunzi Suanjing} is one of the earliest surviving Chinese mathematical texts, written during the 3rd to 5th centuries CE. It was listed as one of the \textit{Ten Computational Canons} during the Tang dynasty.

The treatise consists of three volumes:
\begin{itemize}[leftmargin=1.5em,topsep=2pt,itemsep=1pt]
  \item \textbf{Volume I}: Units of measurement, counting rod notation, arithmetic algorithms
  \item \textbf{Volume II}: Fractions, geometry (areas, volumes), proportional division
  \item \textbf{Volume III}: 36 problems including the famous ``Problem of the Unknown Quantity''—the earliest known statement of the Chinese Remainder Theorem
\end{itemize}
\end{modernbox}

\end{paracol}

\newpage

%% =======================================================================
%% PREFACE
%% =======================================================================
\begin{paracol}{2}

\begin{classicalbox}
\section*{序}
\addcontentsline{toc}{section}{序}

孙子曰：

夫算者，天地之经纬，群生之元首；五常之本末，阴阳之父母；星辰之建号，三光之表里；五行之准平，四时之终始；万物之祖宗，六艺之纲纪。稽群伦之聚散，考二气之降升；推寒暑之迭运，步远近之殊同；观天道精微之兆基，察地理从横之长短；采神祇之所在，极成败之符验；穷道德之理，究性命之情。立规矩，准方圆，谨法度，约尺丈，立权衡，平重轻，剖毫厘，析黍絫；历亿载而不朽，施八极而无疆。散之不可胜究，敛之不盈掌握。向之者富有余，背之者贫且窭；心开者幼冲而即悟，意闭者皓首而难精。夫欲学之者必务量能揆己，志在所专。如是则焉有不成者哉。
\end{classicalbox}

\switchcolumn

\begin{modernbox}
\section*{Preface}

Sun Zi said:

计算（数学），是天地的经纬，万物的起源；是五常（仁义礼智信）的根本，阴阳的父母；是星辰命名的依据，日月星三光的表里；是五行的标准，四季的始终；是万物的祖宗，六艺的纲纪。它可以考察各类事物的聚散，研究阴阳二气的升降；推算寒暑的交替运行，测量远近的异同；观察天道的精微征兆，考察地理的纵横长短；探寻神祇的所在，验证成败的预兆；穷究道德的真理，探究性命的本质。确立规矩，规范方圆，严谨法度，统一尺丈，建立权衡，平衡轻重，剖析毫厘，细数黍粒；历经亿年而不朽，施行八极而无边界。分散开来不可穷尽，收拢起来不满手掌。向它学习的人富富有余，背离它的人贫穷匮乏；心窍开通的人年幼就能领悟，心意闭塞的人白头也难精通。因此想要学习的人必须量力而行、正确评估自己，志向要专一。如果这样，怎么会有学不成的人呢？
\end{modernbox}

\end{paracol}

\newpage

%% =======================================================================
%% VOLUME I
%% =======================================================================
\begin{paracol}{2}

\begin{classicalbox}
\section*{卷上}
\addcontentsline{toc}{section}{卷上}

卷上建立计算的基础：度量衡单位、大数记法、算筹记数法、乘除法则。
\end{classicalbox}

\switchcolumn

\begin{modernbox}
\section*{Volume I}
\addcontentsline{toc}{section}{Volume I (modern)}

Volume I establishes the foundations of calculation: units of measurement, counting rod notation, arithmetic algorithms.
\end{modernbox}

\end{paracol}

\vspace{0.3cm}

\subsubsection*{一、度之所起}
\begin{paracol}{2}
\begin{problemclassical}
度之所起，起于忽。欲知其忽，蚕所生，吐丝为忽。十忽为一秒，十秒为一毫，十毫为一厘，十厘为一分，十分为一寸，十寸为一尺，十尺为一丈，十丈为一引；五十尺为一端；四十尺为一疋；六尺为一步。二百四十步为一亩。三百步为一里。
\end{problemclassical}
\switchcolumn
\begin{problemmodern}
\textbf{【长度单位的起源】}\\
长度单位的起源，从「忽」开始。要知道一忽有多大，就是蚕吐出的丝的粗细。十忽为一秒，十秒为一毫，十毫为一厘，十厘为一分，十分为一寸，十寸为一尺，十尺为一丈，十丈为一引；五十尺为一端；四十尺为一匹；六尺为一步。二百四十步为一亩。三百步为一里。
\end{problemmodern}
\end{paracol}

\vspace{0.3cm}

\subsubsection*{二、称之所起}
\begin{paracol}{2}
\begin{problemclassical}
称之所起，起于黍。十黍为一絫，十絫为一铢，二十四铢为一两，十六两为一斤，三十斤为一钧，四钧为一石。

量之所起，起于粟。六粟为一圭，十圭为一抄，十抄为一撮，十撮为一勺，十勺为一合，十合为一升，十升为一斗，十斗为一斛。斛得六千万粟。
\end{problemclassical}
\switchcolumn
\begin{problemmodern}
\textbf{【重量与容量单位的起源】}\\
重量单位的起源，从「黍」开始。十黍为一絫，十絫为一铢，二十四铢为一两，十六两为一斤，三十斤为一钧，四钧为一石。

容量单位的起源，从「粟」开始。六粟为一圭，十圭为一抄，十抄为一撮，十撮为一勺，十勺为一合，十合为一升，十升为一斗，十斗为一斛。一斛等于六千万粟。
\end{problemmodern}
\end{paracol}

\vspace{0.3cm}

\subsubsection*{三、凡大数之法}
\begin{paracol}{2}
\begin{problemclassical}
凡大数之法，万万曰亿，万亿曰兆，万兆曰京，万京曰陔，万陔曰秭，万秭曰壤，万壤曰沟，万沟曰涧，万涧日正，万万正曰载。
\end{problemclassical}
\switchcolumn
\begin{problemmodern}
\textbf{【大数记法】}\\
大数的记法如下：一万乘一万叫做「亿」；一万乘一亿叫做「兆」；一万乘一兆叫做「京」；一万乘一京叫做「陔」；一万乘一陔叫做「秭」；一万乘一秭叫做「壤」；一万乘一壤叫做「沟」；一万乘一沟叫做「涧」；一万乘一涧叫做「正」；一万乘一正叫做「载」。
\end{problemmodern}
\end{paracol}

\vspace{0.3cm}

\subsubsection*{四、周三径一}
\begin{paracol}{2}
\begin{problemclassical}
周三径一。方五邪七；见邪求方，五之，七而一；见方求邪，七之，五而一。
\end{problemclassical}
\switchcolumn
\begin{problemmodern}
\textbf{【圆周率近似值】}\\
圆周率取周三径一（即 $\pi \approx 3$）。正方形边长为五，则对角线约为七；已知对角线求边长，乘以五再除以七；已知边长求对角线，乘以七再除以五。
\end{problemmodern}
\end{paracol}

\vspace{0.3cm}

\subsubsection*{五、黄金方寸}
\begin{paracol}{2}
\begin{problemclassical}
黄金方寸重一斤。白金方寸重十四两。玉方寸重十二两。铜方寸重七两半。铅方寸重九两半。铁方寸重六两。石方寸重三两。
\end{problemclassical}
\switchcolumn
\begin{problemmodern}
\textbf{【材料比重】}\\
边长为一寸的黄金立方体重一斤。白银一立方寸重十四两。玉一立方寸重十二两。铜一立方寸重七两半。铅一立方寸重九两半。铁一立方寸重六两。石头一立方寸重三两。
\end{problemmodern}
\end{paracol}

\vspace{0.3cm}

\subsubsection*{六、凡算之法}
\begin{paracol}{2}
\begin{problemclassical}
凡算之法，先识其位，一从十横，百立千僵，千十相望，万百相当。
\end{problemclassical}
\switchcolumn
\begin{problemmodern}
\textbf{【算筹记数法】}\\
计算的方法，先要认识数位：个位用纵式，十位用横式，百位用纵式（直立），千位用横式（平放）。千位和十位的记法相互对应，万位和百位的记法相互对应。

这描述了中国传统的算筹记数法：
\begin{itemize}[leftmargin=1.5em,topsep=1pt,itemsep=0pt]
  \item 个位（1--9）用\textbf{纵式}
  \item 十位（10--90）用\textbf{横式}
  \item 百位（100--900）用\textbf{纵式}
  \item 千位（1000--9000）用\textbf{横式}
\end{itemize}
每个数位的纵式与横式交替排列。
\end{problemmodern}
\end{paracol}

\vspace{0.3cm}

\subsubsection*{七、凡乘之法}
\begin{paracol}{2}
\begin{problemclassical}
凡乘之法，重置其位。上下相观，上位有十步至十，有百步至百，有千步至千。以上命下，所得之数列于中位。言十即过，不满自如。上位乘讫者先去之。下位乘讫者则俱退之。六不积，五不只。上下相乘，至尽则已。
\end{problemclassical}
\switchcolumn
\begin{problemmodern}
\textbf{【乘法法则】}\\
做乘法时，重新排列被乘数和乘数的位置。上下对齐观察：如果被乘数有十就进到十位，有百就进到百位，有千就进到千位。用上面的数字去乘下面的数字，所得的结果写在中位。满十就进位，不满十就保留原位。上位的数字乘完就撤去。下位的数字乘完就都退后一位。算筹中六不用五上加一表示，五不单独表示。上下相乘，直到全部乘完为止。
\end{problemmodern}
\end{paracol}

\vspace{0.3cm}

\subsubsection*{八、凡除之法}
\begin{paracol}{2}
\begin{problemclassical}
凡除之法，与乘正异。乘得在中央，除得在上方。假令六为法，百为实。以六除百，当进之二等，令在正百下，以六除一，则法多而实少，不可除，故当退就十位。以法除实，言一六而折百为四十，故可除。若实多法少，自当百之，不当复退。故或步法十者置于十位，百者置于百位。上位有空绝者，法退二位。余法皆如乘时。实有余者，以法命之，以法为母。实余为子。
\end{problemclassical}
\switchcolumn
\begin{problemmodern}
\textbf{【除法法则】}\\
除法的方法和乘法正好相反。乘法的结果在中央，除法的结果在上方。假设六作为除数（法），一百作为被除数（实）。用六去除一百，应当进两位，放在百位下面。用六去除一，除数比被除数大，无法除，所以应当退到十位。用除数去除被除数：说「一六」并从一百中减去六十，余四十，所以可以除。如果被除数大除数小，自然应该放在百位，不应再退。所以如果除数是十就放在十位，是百就放在百位。上位如果为空，除数就退两位。其余方法和乘法一样。被除数如果有余数，就用除数来命名分数，以除数为分母，余数为分子。
\end{problemmodern}
\end{paracol}

\newpage

\subsubsection*{九、以粟求糲米}
\begin{paracol}{2}
\begin{problemclassical}
以粟求糲米，三之，五而一。以糲米求粟，五之，三而一。以糲米求饭，五之，二而一。以粟米求糲饭，六之，四而一。以糲饭求糲米，二之，五而一。以綹米求饭，八之，四而一。
\end{problemclassical}
\switchcolumn
\begin{problemmodern}
\textbf{【粮谷换算】}\\
粟米换算成糙米，乘以三再除以五。糙米换算回粟米，乘以五再除以三。糙米煮成饭，乘以五再除以二。粟米换算成糙米饭，乘以六再除以四。糙米饭还原成糙米，乘以二再除以五。綹米煮成饭，乘以八再除以四。
\end{problemmodern}
\end{paracol}

\vspace{0.3cm}

\subsubsection*{一〇、十分减一者}
\begin{paracol}{2}
\begin{problemclassical}
十分减一者，以二乘，二十除。减二者，以四乘，二十除。减三者，以六乘，二十除。减四者，以八乘，二十除。减五者，以十乘，二十除。减六者，以十二乘，二十除。减七者，以十四乘，二十除。减八者，以十六乘，二十除。减九者，以十八乘，二十除。
\end{problemclassical}
\switchcolumn
\begin{problemmodern}
\textbf{【十分之一折扣】}\\
减十分之一的，乘以二再除以二十。减十分之二的，乘以四再除以二十。减十分之三的，乘以六再除以二十。减十分之四的，乘以八再除以二十。减十分之五的，乘以十再除以二十。减十分之六的，乘以十二再除以二十。减十分之七的，乘以十四再除以二十。减十分之八的，乘以十六再除以二十。减十分之九的，乘以十八再除以二十。
\end{problemmodern}
\end{paracol}

\vspace{0.3cm}

\subsubsection*{一六、九九八十一}
\begin{paracol}{2}
\begin{problemclassical}
\textbf{问：}九九八十一，自相乘，得几何？

\textbf{答曰：}六千五百六十一。

\textbf{术曰：}重置其位，以上八呼下八，八八六十四，即下六千四百于中位。以上八呼下一，一八如八，即于中位下八十。退下位一等，收上位八十。以上位一呼下八，一八如八，即于中位下八十。以上一呼下一，一一如一，即于中位下一。上下位俱收，中位即得六千五百六十一。
\end{problemclassical}
\switchcolumn
\begin{problemmodern}
\textbf{【题目】} 81乘以81，结果是多少？

\textbf{答案：} 6561。

\textbf{解法：} 重新排列算筹位置，用上排的8去乘下排的8，八八六十四，在中位写下6400。用上排的8去乘下排的1，一八得八，在中位写下80。把下位退一位，收去上位的80。用上位的1去乘下位的8，一八得八，在中位写下80。用上位的1去乘下位的1，一一得一，在中位写下1。上下位都收去后，中位得到6561。
\end{problemmodern}
\end{paracol}

\vspace{0.3cm}

\subsubsection*{一七、六千五百六十一}
\begin{paracol}{2}
\begin{problemclassical}
\textbf{问：}六千五百六十一，九人分之，问人得几何？

\textbf{答曰：}七百二十九。

\textbf{术曰：}先置六千五百六十一于中位，为实。下列九人为法。上位置七百，以上七呼下九，七九六十三，即除中位六千三百。退下位一等，即上位置二十。以上二呼下九，二九十八，即除中位一百八十。又更退下位一等，即上位更置九，即以上九呼下九，九九八十一，即除中位八十一。中位尽，收下位。上位所得即人之所得。自八八六十四至一一如一，并准此。
\end{problemclassical}
\switchcolumn
\begin{problemmodern}
\textbf{【题目】} 6561由九个人平分，每人分得多少？

\textbf{答案：} 729。

\textbf{解法：} 先把6561放在中位作为被除数（实）。下面放9人作为除数（法）。上位置700，用上面的7乘下面的9，七九六十三，从中位减去6300。下位退一位，上位置20。用上面的2乘下面的9，二九十八，从中位减去180。再退下位一位，上位再放9，用上面的9乘下面的9，九九八十一，从中位减去81。中位减尽，收下位。上位所得就是每人分得的数量。从八八六十四到一一如一，都照这个方法。
\end{problemmodern}
\end{paracol}

\newpage

%% =======================================================================
%% VOLUME II
%% =======================================================================
\begin{paracol}{2}

\begin{classicalbox}
\section*{卷中}
\addcontentsline{toc}{section}{卷中}

卷中共二十八题，涉及分数运算、面积体积、比例分配等。
\end{classicalbox}

\switchcolumn

\begin{modernbox}
\section*{Volume II}
\addcontentsline{toc}{section}{Volume II (modern)}

Volume II contains 28 problems dealing with fractions, areas, volumes, proportional distribution, and practical calculations.
\end{modernbox}

\end{paracol}

\vspace{0.3cm}

\subsubsection*{卷中第一题：约分}
\begin{paracol}{2}
\begin{problemclassical}
\textbf{问：}今有一十八分之一十二。问约之得几何？

\textbf{答曰：}三分之二。

\textbf{术曰：}置十八分在下，一十二分在上。副置二位，以少减多，等数得六，为法。约之，即得。
\end{problemclassical}
\switchcolumn
\begin{problemmodern}
\textbf{【题目】} 现在有12/18。问约分后等于多少？

\textbf{答案：} 2/3。

\textbf{解法：} 把分母18放在下面，分子12放在上面。另放两个数在旁边，用较大数减去较小数（辗转相减），等数（最大公约数）得到6，作为除数。分子分母都除以6，就得到答案。
\end{problemmodern}
\end{paracol}

\vspace{0.3cm}

\subsubsection*{卷中第二题：合分}
\begin{paracol}{2}
\begin{problemclassical}
\textbf{问：}今有三分之一，五分之二。问合之得几何？

\textbf{答曰：}一十五分之一十一。

\textbf{术曰：}置三分、五分在右方，之一、之二在左方。母互乘子，五分之二得六，三分之一得五。并之，得一十一，为实。右方二母相乘，得一十五，为法。不满法，以法命之，即得。
\end{problemclassical}
\switchcolumn
\begin{problemmodern}
\textbf{【题目】} 现在有1/3和2/5。问相加等于多少？

\textbf{答案：} 11/15。

\textbf{解法：} 把分母3和5放在右边，分子1和2放在左边。交叉相乘：2/5的分子变成6（$2 \times 3$），1/3的分子变成5（$1 \times 5$）。相加得到11，作为被除数（分子）。两个分母相乘得到15，作为除数（分母）。11不足15，所以答案就是11/15。
\end{problemmodern}
\end{paracol}

\vspace{0.3cm}

\subsubsection*{卷中第三题：减分}
\begin{paracol}{2}
\begin{problemclassical}
\textbf{问：}今有九分之八，减其五分之一。问余几何？

\textbf{答曰：}四十五分之三十一。

\textbf{术曰：}置九分、五分在右方，之八、之一在左方。母互乘子，五分之一得九，九分之八得四十。以少减多，余三十一，为实。母相乘得四十五，为法。不满法，以法命之，即得。
\end{problemclassical}
\switchcolumn
\begin{problemmodern}
\textbf{【题目】} 现在有8/9，减去它的1/5。问还剩多少？

\textbf{答案：} 31/45。

\textbf{解法：} 把分母9和5放在右边，分子8和1放在左边。交叉相乘：1/5的分子变成9（$1 \times 9$），8/9的分子变成40（$8 \times 5$）。用大数减小数，余31，作为分子。分母相乘得45，作为分母。31不足45，所以答案是31/45。
\end{problemmodern}
\end{paracol}

\vspace{0.3cm}

\subsubsection*{卷中第四题：平分}
\begin{paracol}{2}
\begin{problemclassical}
\textbf{问：}今有三分之一，三分之二，四分之三。问减多益少，几何而平？

\textbf{答曰：}减四分之三者二，减三分之二者一，并以益三分之一，而各平于一十二分之七。
\end{problemclassical}
\switchcolumn
\begin{problemmodern}
\textbf{【题目】} 现在有1/3、2/3、3/4这三个分数。问：从多的里面减去一些，补给少的，使得三个数相等，应该怎么分？

\textbf{答案：} 从3/4中减去2/12，从2/3中减去1/12，都补给1/3，这样每个数都等于7/12。
\end{problemmodern}
\end{paracol}

\vspace{0.3cm}

\subsubsection*{卷中第五题：粟求糯米}
\begin{paracol}{2}
\begin{problemclassical}
\textbf{问：}今有粟一斗，问为糯米几何？

\textbf{答曰：}六升。

\textbf{术曰：}置粟一斗，十升。以糲米率三十乘之，得三百升，为实。以粟率五十为法，除之，即得。
\end{problemclassical}
\switchcolumn
\begin{problemmodern}
\textbf{【题目】} 现在有粟米一斗，问能换糯米多少？

\textbf{答案：} 六升。

\textbf{解法：} 把一斗粟米换算成十升。用糯米率30去乘，得到300升，作为被除数。用粟率50作为除数去除，即得6升。
\end{problemmodern}
\end{paracol}

\vspace{0.3cm}

\subsubsection*{卷中第九题：屋基用砖}
\begin{paracol}{2}
\begin{problemclassical}
\textbf{问：}今有屋基南北三丈，东西六丈，欲以砖砌之。凡积二尺，用砖五枚。问计几何？

\textbf{答曰：}四千五百枚。

\textbf{术曰：}置东西六丈，以南北三丈乘之，得一千八百尺。以五乘之，得九千尺。以二除之，即得。
\end{problemclassical}
\switchcolumn
\begin{problemmodern}
\textbf{【题目】} 有一座房屋的基座，南北长三丈，东西宽六丈，想用砖砌成。每二立方尺用五块砖。问一共需要多少块砖？

\textbf{答案：} 4500块。

\textbf{解法：} 东西六丈乘以南北三丈，得到1800平方尺。乘以5得到9000，除以2，即得4500块。
\end{problemmodern}
\end{paracol}

\newpage

\subsubsection*{卷中第一〇题：圆窖受粟}
\begin{paracol}{2}
\begin{problemclassical}
\textbf{问：}今有圆窖下周二百八十六尺，深三丈六尺。问受粟几何？

\textbf{答曰：}一十五万一千四百七十四斛七升二十七分升之十一。

\textbf{术曰：}置周二百八十六尺，自相乘，得八万一千七百九十六尺。以深三丈六尺乘之，得二百九十四万四千六百五十六尺。以十二除之，得二十四万五千三百八十八尺。以斛法一尺六寸二分除之，即得。
\end{problemclassical}
\switchcolumn
\begin{problemmodern}
\textbf{【题目】} 有一个圆形地窖，底面周长286尺，深36尺。问能装多少粟米？

\textbf{答案：} 151474斛7升又11/27升。

\textbf{解法：} 周长286尺自乘，得81796平方尺。乘以深度36尺，得2944656立方尺。除以12，得245388立方尺。再除以每斛容积1.62立方尺，即得。
\end{problemmodern}
\end{paracol}

\vspace{0.3cm}

\subsubsection*{卷中第一三题：圆田面积}
\begin{paracol}{2}
\begin{problemclassical}
\textbf{问：}今有圆田周三百步，径一百步。问得田几何？

\textbf{答曰：}三十一亩奇六十步。

\textbf{术曰：}先置周三百步，半之，得一百五十步。又置径一百步，半之，得五十步。相乘，得七千五百步。以亩法二百四十步除之，即得。
\end{problemclassical}
\switchcolumn
\begin{problemmodern}
\textbf{【题目】} 有一块圆形田地，周长300步，直径100步。问田地面积是多少？

\textbf{答案：} 31亩余60步。

\textbf{解法：} 把周长300步取半，得150步。再把直径100步取半，得50步。两个半数相乘，得7500平方步。用每亩240步去除，即得31亩余60步。
\end{problemmodern}
\end{paracol}

\vspace{0.3cm}

\subsubsection*{卷中第一六题：堤的体积}
\begin{paracol}{2}
\begin{problemclassical}
\textbf{问：}今有堤，下广五丈，上广三丈，高二丈，长六十尺。欲以一千尺为一方，问计几何？

\textbf{答曰：}四十八方。

\textbf{术曰：}置堤上广三丈，下广五丈。并之，得八丈。半之，得四丈。以高二丈乘之，得八百尺。以长六十尺乘之，得四万八千。以一千尺除之，即得。
\end{problemclassical}
\switchcolumn
\begin{problemmodern}
\textbf{【题目】} 有一道堤坝，下底宽5丈，上顶宽3丈，高2丈，长60尺。以1000立方尺为一个计量单位，问共有多少个单位？

\textbf{答案：} 48方。

\textbf{解法：} 上宽3丈加下宽5丈，得8丈。取半得4丈。乘以高2丈得800平方尺。再乘以长60尺得48000立方尺。除以1000，即得48方。
\end{problemmodern}
\end{paracol}

\vspace{0.3cm}

\subsubsection*{卷中第一八题：开平方}
\begin{paracol}{2}
\begin{problemclassical}
\textbf{问：}今有积二十三万四千五百六十七步。问为方几何？

\textbf{答曰：}四百八十四步九百六十八分步之三百一十一。

\textbf{术曰：}置积二十三万四千五百六十七步，为实。次借一算为下法。步之，超一位，至百而止。商置四百……上商得四百八十四，下法得九百六十八，不尽三百一十一。是为方四百八十四步九百六十八分步之三百一十一。
\end{problemclassical}
\switchcolumn
\begin{problemmodern}
\textbf{【题目】} 有一个面积为234567平方步的正方形。问它的边长是多少？

\textbf{答案：} $484\,\frac{311}{968}$ 步。

\textbf{解法：} 这是中国古代的「开平方」算法——世界上最早的十进制开平方方法。把234567作为被开方数（实）。借一算筹作为下法。隔一位跳一步，到百位为止。商上放400……最终商得484，下法得968，余数311。所以边长为484又311/968步。
\end{problemmodern}
\end{paracol}

\newpage

%% =======================================================================
%% VOLUME III
%% =======================================================================
\begin{paracol}{2}

\begin{classicalbox}
\section*{卷下}
\addcontentsline{toc}{section}{卷下}

卷下共三十六题，涵盖各种实际应用问题，包括粮谷分配、土地测量，以及著名的「物不知数」题——中国剩余定理的最早记载。
\end{classicalbox}

\switchcolumn

\begin{modernbox}
\section*{Volume III}
\addcontentsline{toc}{section}{Volume III (modern)}

Volume III contains 36 problems covering a wide range of practical applications, from grain distribution and land measurement to the famous ``Problem of the Unknown Quantity''—the earliest known statement of the Chinese Remainder Theorem.
\end{modernbox}
\end{paracol}

\vspace{0.3cm}

\subsubsection*{卷下第五题：围棋盘}
\begin{paracol}{2}
\begin{problemclassical}
\textbf{问：}今有棋局方一十九道。问用棋几何？

\textbf{答曰：}三百六十一。

\textbf{术曰：}置一十九道，自相乘之，即得。
\end{problemclassical}
\switchcolumn
\begin{problemmodern}
\textbf{【题目】} 有一个围棋盘，每边有19道线。问棋盘上有多少个交叉点（用多少颗棋子）？

\textbf{答案：} 361个。

\textbf{解法：} 19道线自乘，即 $19 \times 19 = 361$ 个交叉点。
\end{problemmodern}
\end{paracol}

\vspace{0.3cm}

\subsubsection*{卷下第一八题：立表测影}
\begin{paracol}{2}
\begin{problemclassical}
\textbf{问：}今有竿不知长短，度其影得一丈五尺。别立一表，长一尺五寸，影得五寸。问竿长几何？

\textbf{答曰：}四丈五尺。

\textbf{术曰：}置竿影一丈五尺，以表长一尺五寸乘之，上十之，得二十二丈五尺。以表影五寸除之，即得。
\end{problemclassical}
\switchcolumn
\begin{problemmodern}
\textbf{【题目】} 有一根竹竿，不知道有多长。量它的影子，长1丈5尺。另外立一根标竿，长1尺5寸，影子长5寸。问竹竿长多少？

\textbf{答案：} 4丈5尺。

\textbf{解法：} 这是中国古代利用相似三角形进行比例计算的例子。竹竿影长1.5丈乘以标竿长1.5尺，再乘以10，得22.5丈。除以标竿影长0.5尺，即得4.5丈（45尺）。
\end{problemmodern}
\end{paracol}

\vspace{0.3cm}

\subsubsection*{卷下第一九题：器中余米}
\begin{paracol}{2}
\begin{problemclassical}
\textbf{问：}今有器中米，不知其数。前人取半，中人三分取一，后人四分取一，余米一斗五升。问本米几何？

\textbf{答曰：}六斗。

\textbf{术曰：}置余米一斗五升，以六乘之，得九斗。以二除之，得四斗五升。以四乘之，得一斛八斗。以三除之，即得。
\end{problemclassical}
\switchcolumn
\begin{problemmodern}
\textbf{【题目】} 容器中有一些米，不知道有多少。第一个人取走一半，第二个人取走剩余的三分之一，第三个人取走再剩余的四分之一，最后还剩1斗5升。问原来有多少米？

\textbf{答案：} 六斗。

\textbf{解法：} 把剩余的1斗5升乘以6，得9斗。除以2得4.5斗。乘以4得1斛8斗。除以3，即得6斗。
\end{problemmodern}
\end{paracol}

\vspace{0.3cm}

\subsubsection*{卷下第二〇题：人车间题}
\begin{paracol}{2}
\begin{problemclassical}
\textbf{问：}今有三人共车，二车空；二人共车，九人步。问人与车各几何？

\textbf{答曰：}一十五车。三十九人。

\textbf{术曰：}置二车，以三乘之，得六。加步者九人，得车一十五。欲知人者，以二乘车，加九人，即得。
\end{problemclassical}
\switchcolumn
\begin{problemmodern}
\textbf{【题目】} 现在有一些人坐车：如果每3人坐一辆车，会空出2辆车；如果每2人坐一辆车，会有9人需要步行。问一共有多少人和多少辆车？

\textbf{答案：} 15辆车，39人。

\textbf{解法：} 把空出的2辆车乘以3（每车3人），得6。加上步行的9人，得15辆车。要求人数：用2人乘以15辆车，加9人，即得39人。
\end{problemmodern}
\end{paracol}

\vspace{0.3cm}

\subsubsection*{卷下第二一题：河上荡杯}
\begin{paracol}{2}
\begin{problemclassical}
\textbf{问：}今有妇人河上荡杯。津吏问曰：「杯何以多？」妇人曰：「家有客。」津吏曰：「客几何？」妇人曰：「二人共饭，三人共羹，四人共肉，凡用杯六十五。不知客几何？」

\textbf{答曰：}六十人。

\textbf{术曰：}置六十五杯，以一十二乘之，得七百八十，以十三除之，即得。
\end{problemclassical}
\switchcolumn
\begin{problemmodern}
\textbf{【题目】} 有一位妇女在河边洗杯子。渡口官吏问：「为什么有这么多杯子？」妇女说：「家里有客人。」官吏问：「有多少客人？」妇女说：「两个人共用一碗饭，三个人共用一碗汤，四个人共用一碗肉，一共用了65个碗。不知道有多少客人？」

\textbf{答案：} 60人。

\textbf{解法：} 把65个碗乘以12（2、3、4的最小公倍数的约数），得780，再除以13（$6+4+3=13$），即得60人。
\end{problemmodern}
\end{paracol}

\newpage

%% =======================================================================
%% THE CHINESE REMAINDER THEOREM
%% =======================================================================
\begin{paracol}{2}

\begin{tcolorbox}[colback=crtgold!8,colframe=crtgold,boxrule=1.5pt,
  left=10pt,right=10pt,top=8pt,bottom=8pt,arc=3mm]
\section*{\color{crtgold}物不知数}
\addcontentsline{toc}{section}{物不知数 — 中国剩余定理}

\begin{center}
\textbf{\large 卷下第26问 — 中国剩余定理}
\end{center}

\vspace{0.2cm}

\textbf{问：}今有物，不知其数。三三数之，剩二；五五数之，剩三；七七数之，剩二。问物几何？

\vspace{0.2cm}

\textbf{答曰：}二十三。

\vspace{0.2cm}

\textbf{术曰：}「三三数之，剩二」，置一百四十；「五五数之，剩三」，置六十三；「七七数之，剩二」，置三十。并之，得二百三十三。以二百一十减之，即得。

凡三三数之，剩一，则置七十；五五数之，剩一，则置二十一；七七数之，剩一，则置十五。一百六以上，以一百五减之，即得。
\end{tcolorbox}

\switchcolumn

\begin{tcolorbox}[colback=crtgold!8,colframe=crtgold,boxrule=1.5pt,
  left=10pt,right=10pt,top=8pt,bottom=8pt,arc=3mm]
\section*{\color{crtgold}Problem of the Unknown Quantity}
\addcontentsline{toc}{section}{物不知数 (modern)}

\begin{center}
\textbf{\large Problem 26, Volume III — The Chinese Remainder Theorem}
\end{center}

\vspace{0.2cm}

\textbf{【题目】} 有一些物品，不知道有多少个。三个三个地数，最后剩两个；五个五个地数，最后剩三个；七个七个地数，最后剩两个。问这些物品一共有多少个？

\vspace{0.2cm}

\textbf{【答案】} 23个。

\vspace{0.2cm}

\textbf{【解法】} 「三个三个数剩二」，先记下140；「五个五个数剩三」，记下63；「七个七个数剩二」，记下30。把这三个数相加，得233。减去210（$3 \times 5 \times 7 = 105$ 的两倍），即得23。

一般规律：三个三个数剩一，就记70；五个五个数剩一，就记21；七个七个数剩一，就记15。总数超过106，就减去105，即得答案。
\end{tcolorbox}

\end{paracol}

\vspace{0.5cm}

\subsubsection*{数学分析}
\begin{paracol}{2}
\begin{classicalbox}
\textbf{孙子歌诀：}

三人同行七十稀，五树梅花廿一枝，七子团圆正半月，除百零五便得知。

\begin{itemize}[leftmargin=1.5em,topsep=2pt,itemsep=1pt]
  \item 除以3的余数 $\times$ 70
  \item 除以5的余数 $\times$ 21
  \item 除以7的余数 $\times$ 15
  \item 总和超过105就减去105
\end{itemize}

\textbf{验证：} $2 \times 70 + 3 \times 21 + 2 \times 15 = 140 + 63 + 30 = 233$

$233 - 2 \times 105 = 23$
\end{classicalbox}
\switchcolumn
\begin{modernbox}
\textbf{Modern Mathematical Analysis:}

The problem asks us to find the smallest positive integer $N$ satisfying:
\begin{align*}
N &\equiv 2 \pmod{3} \\
N &\equiv 3 \pmod{5} \\
N &\equiv 2 \pmod{7}
\end{align*}

The key numbers 70, 21, and 15 have special properties:
\begin{itemize}[leftmargin=1.5em,topsep=2pt,itemsep=1pt]
  \item $70 \equiv 1 \pmod{3}$, $70 \equiv 0 \pmod{5,7}$
  \item $21 \equiv 1 \pmod{5}$, $21 \equiv 0 \pmod{3,7}$
  \item $15 \equiv 1 \pmod{7}$, $15 \equiv 0 \pmod{3,5}$
\end{itemize}

General formula: $N = 70r_1 + 21r_2 + 15r_3 - 105k$

\textbf{Check:} $23 \div 3 = 7$ r.2\ \checkmark,\ $23 \div 5 = 4$ r.3\ \checkmark,\ $23 \div 7 = 3$ r.2\ \checkmark
\end{modernbox}
\end{paracol}

\newpage

%% =======================================================================
%% REMAINING PROBLEMS FROM VOLUME III
%% =======================================================================

\subsubsection*{卷下第二七题：鸟兽同笼}
\begin{paracol}{2}
\begin{problemclassical}
\textbf{问：}今有兽六首四足，禽四首二足。上有七十六首，下有四十六足。问禽、兽各几何？

\textbf{答曰：}八兽，七禽。

\textbf{术曰：}倍足以减首，余，半之，即禽。四乘禽数，减足，余，半之，即兽。
\end{problemclassical}
\switchcolumn
\begin{problemmodern}
\textbf{【题目】} 现在有一些怪兽和禽鸟：怪兽有6个头4只脚，禽鸟有4个头2只脚。从上面数共有76个头，从下面数共有46只脚。问禽鸟和怪兽各有多少只？

\textbf{答案：} 怪兽8只，禽鸟7只。

\textbf{解法：} 把脚数加倍再减去头数，余数取半就是禽鸟数。用禽鸟数乘以4，从脚数中减去，余数取半就是怪兽数。
\end{problemmodern}
\end{paracol}

\vspace{0.3cm}

\subsubsection*{卷下第二九题：百鹿入城}
\begin{paracol}{2}
\begin{problemclassical}
\textbf{问：}今有百鹿入城，家取一鹿，不尽；又三家共一鹿，适尽。问城中家几何？

\textbf{答曰：}七十五家。

\textbf{术曰：}以盈不足取之。假令七十二家，鹿不尽四。令之九十家，鹿不足二十……
\end{problemclassical}
\switchcolumn
\begin{problemmodern}
\textbf{【题目】} 有100头鹿进城，如果每家分一头鹿，分不完；又如果三家共分一头鹿，正好分完。问城里有多少户人家？

\textbf{答案：} 75家。

\textbf{解法：} 用盈不足术求解。假设72家，鹿多余4头。假设90家，鹿不足20头。用盈不足公式即可求得75家。

这是中国古代的「盈不足术」——一种解线性方程的巧妙方法。
\end{problemmodern}
\end{paracol}

\vspace{0.3cm}

\subsubsection*{卷下第三一题：雉兔同笼}
\begin{paracol}{2}
\begin{problemclassical}
\textbf{问：}今有雉、兔同笼，上有三十五头，下九十四足。问雉、兔各几何？

\textbf{答曰：}雉二十三。兔一十二。

\textbf{术曰：}上置三十五头，下置九十四足。半其足，得四十七。以少减多，再命之，上三除下三，上五除下五。下有一除上一，下有二除上二，即得。
\end{problemclassical}
\switchcolumn
\begin{problemmodern}
\textbf{【题目】} 现在有一些野鸡和兔子关在同一个笼子里，从上面数有35个头，从下面数有94只脚。问野鸡和兔子各有多少只？

\textbf{答案：} 野鸡23只，兔子12只。

\textbf{解法：} 这是中国古代著名的「鸡兔同笼」问题。把上面的头数35和下面的脚数94列出来。脚数取半得47。用较大数减较小数，再进行操作，即可分别求出野鸡和兔子的数量。

现代解法：设野鸡$x$只，兔子$y$只，则 $x+y=35$，$2x+4y=94$，解得 $x=23$，$y=12$。
\end{problemmodern}
\end{paracol}

\vspace{0.3cm}

\subsubsection*{卷下第三三题：洛阳至长安}
\begin{paracol}{2}
\begin{problemclassical}
\textbf{问：}今有长安、洛阳相去九百里。车轮一匝一丈八尺。欲自洛阳至长安，问轮匝几何？

\textbf{答曰：}九万匝。

\textbf{术曰：}置九百里，以三百步乘之，得二十七万步。又以六尺乘之，得一百六十二万尺。以车轮一丈八尺为法。除之，即得。
\end{problemclassical}
\switchcolumn
\begin{problemmodern}
\textbf{【题目】} 长安和洛阳相距900里。车轮转一圈走1丈8尺。想从洛阳到长安，问车轮要转多少圈？

\textbf{答案：} 90000圈。

\textbf{解法：} 把900里换算成步：$900 \times 300 = 270000$步。再换算成尺：$270000 \times 6 = 1620000$尺。用车轮周长18尺去除，即得90000圈。
\end{problemmodern}
\end{paracol}

\vspace{0.3cm}

\subsubsection*{卷下第三四题：出门望九堤}
\begin{paracol}{2}
\begin{problemclassical}
\textbf{问：}今有出门望见九堤。堤有九木，木有九枝，枝有九巢，巢有九禽，禽有九雏，雏有九毛，毛有九色。问各几何？

\textbf{答曰：}木八十一。枝七百二十九。巢六千五百六十一。禽五万九千四十九。雏五十三万一千四百四十一。毛四百七十八万二千九百六十九。色四千三百四万六千七百二十一。

\textbf{术曰：}置九堤，以九乘之，得木之数。又以九乘之，得枝之数……又以九乘之，得色之数。
\end{problemclassical}
\switchcolumn
\begin{problemmodern}
\textbf{【题目】} 有一个人出门看见9道堤坝。每道堤坝有9棵树，每棵树有9根枝，每根枝有9个巢，每个巢有9只鸟，每只鸟有9只雏，每只雏有9根毛，每根毛有9种颜色。问各有多少？

\textbf{答案：} 树81棵，枝729根，巢6561个，鸟59049只，雏531441只，毛4782969根，颜色43046721种。

\textbf{解法：} 这是一道连乘题，每次乘以9：

\begin{tabular}{@{}ll@{}}
树 & $9 \times 9 = 81$ \\
枝 & $81 \times 9 = 729$ \\
巢 & $729 \times 9 = 6561$ \\
鸟 & $6561 \times 9 = 59049$ \\
雏 & $59049 \times 9 = 531441$ \\
毛 & $531441 \times 9 = 4782969$ \\
色 & $4782969 \times 9 = 43046721$ \\
\end{tabular}
\end{problemmodern}
\end{paracol}

\newpage

\subsubsection*{卷下第三五题：三女归家}
\begin{paracol}{2}
\begin{problemclassical}
\textbf{问：}今有三女，长女五日一归，中女四日一归，少女三日一归。问三女几何日相会？

\textbf{答曰：}六十日。

\textbf{术曰：}置长女五日、中女四日、少女三日，于右方。各列一算于左方。维乘之，各得所到数：长女十二到，中女十五到，少女二十到。又各以归日乘到数，即得。
\end{problemclassical}
\switchcolumn
\begin{problemmodern}
\textbf{【题目】} 有三姐妹，大姐每5天回一次娘家，二姐每4天回一次，小妹每3天回一次。问三姐妹最少多少天能在娘家相会？

\textbf{答案：} 60天。

\textbf{解法：} 这是中国古代的「最小公倍数」问题。把大姐5天、二姐4天、小妹3天放在右边。各列一个算筹在左边。交叉相乘：大姐回娘家12次，二姐15次，小妹20次。再各自用归日乘以到数，都得到60。

即求3、4、5的最小公倍数：$\text{lcm}(3,4,5) = 60$。
\end{problemmodern}
\end{paracol}

\vspace{0.3cm}

\subsubsection*{卷下第三六题：孕妇推男女}
\begin{paracol}{2}
\begin{problemclassical}
\textbf{问：}今有孕妇行年二十九，难九月。未知所生？

\textbf{答曰：}生男。

\textbf{术曰：}置四十九，加难月，减行年。所余，以天除一，地除二，人除三，四时除四，五行除五，六律除六，七星除七，八风除八，九州除九。其不尽者，奇则为男，耦则为女。
\end{problemclassical}
\switchcolumn
\begin{problemmodern}
\textbf{【题目】} 有一位孕妇，今年29岁，怀孕9个月。问生男孩还是女孩？

\textbf{答案：} 生男孩。

\textbf{解法：} 把49加上怀孕的月数（9），再减去孕妇的年龄（29），得到余数。用余数依次减去：天一、地二、人三、四时四、五行五、六律六、七星七、八风八、九州九。如果最后余数是奇数就生男孩，偶数就生女孩。

$49 + 9 - 29 = 29$，然后依次减 $1+2+3+4+5+6+7 = 28$，余1（奇数），所以预测生男孩。

\textit{注：这是古代的迷信算法，没有科学依据。}
\end{problemmodern}
\end{paracol}

\newpage

%% =======================================================================
%% EPILOGUE
%% =======================================================================
\begin{paracol}{2}

\begin{classicalbox}
\section*{结语}
\addcontentsline{toc}{section}{结语}

《孙子算经》是古代中国数学的一部杰出著作。虽然书中涵盖的多是基础算术——度量衡、九九乘法表、简单分数，但其最伟大的遗产在于那道简短的「物不知数」题，它为现代数论打开了一扇大门。

中国剩余定理，这一方法在西方后来被称为「中国剩余定理」，在欧洲要再过一千多年才被完整形式化。孙子书中描述的算法包含了现代定理的核心要素：构造具有特定整除性质的辅助数，将它们与给定的余数组合，再对最小公倍数取模。

这部著作的影响远远超出了中国。这道题作为数学趣题在整个东亚流传，而秦九韶发展的一般方法代表了中世纪世界数学的伟大成就之一。
\end{classicalbox}

\switchcolumn

\begin{modernbox}
\section*{Epilogue}

The \textit{Sunzi Suanjing} stands as a remarkable testament to the sophistication of ancient Chinese mathematics. While the text covers basic arithmetic—units of measurement, multiplication tables, and simple fractions—its greatest legacy lies in the brief but brilliant ``Problem of the Unknown Quantity'' that opens the door to modern number theory.

The Chinese Remainder Theorem, as this method came to be known in the West, would not be fully formalized in Europe for over a millennium. The algorithm described by Sunzi contains the essential ingredients of the modern theorem: the construction of auxiliary numbers with specific divisibility properties, their combination with the given remainders, and reduction modulo the least common multiple.

The influence of this text extended far beyond China. The problem circulated as a mathematical puzzle throughout East Asia, and the general method developed by Qin Jiushao represents one of the great achievements of medieval mathematics anywhere in the world.
\end{modernbox}

\end{paracol}

\vspace{1cm}
\begin{center}
\longline
\end{center}

\vspace{0.5cm}

\begin{center}
\textit{据中华经典古籍库《孙子算经》原文排版}\\
\textit{（原载于中华古籍资源库 / ctext.org）}
\end{center}

\end{document}

